Ce travail présente la conception et la réalisation d'un \textbf{Système de Détection d'Anomalies Comportementales Personnelles (PBADS)} basé sur une architecture microservices. Le système analyse les données quotidiennes des utilisateurs (sommeil, activité physique, alimentation, humeur) et détecte en temps réel les anomalies comportementales à l'aide d'algorithmes d'apprentissage automatique non supervisés.

La solution repose sur une \textbf{architecture microservices} avec 13 services distincts, incluant un service d'authentification (JWT), un service de données, un service de modèles ML, un service d'inférence pour la détection d'anomalies, un service d'alertes, et un service de recommandations alimenté par Google Gemini AI. La communication entre services utilise Kafka pour la messagerie asynchrone et HTTP/REST pour les appels synchrones, avec un mécanisme de fallback pour garantir la résilience.

Le modèle de données structure les entités principales : \texttt{User}, \texttt{DailyData}, \texttt{Alert} et \texttt{Recommendation}. Le système utilise des algorithmes de détection d'anomalies (Isolation Forest et One-Class SVM) pour identifier les écarts par rapport aux comportements normaux, génère des alertes avec des niveaux de sévérité (NORMAL, LOW, MEDIUM, HIGH), et fournit des recommandations personnalisées via l'intégration avec Google Gemini AI.

Les bénéfices observés : \textbf{détection automatique des anomalies}, \textbf{recommandations personnalisées}, \textbf{notifications en temps réel}, et \textbf{visualisation interactive} des données comportementales via un tableau de bord avec filtres temporels. Le système offre une scalabilité élevée grâce à l'architecture microservices et une tolérance aux pannes via Resilience4j.

\textbf{Mots-clés :} Détection d'anomalies, Machine Learning, Microservices, Spring Boot, Kafka, PostgreSQL, Isolation Forest, One-Class SVM, Google Gemini AI, Recommandations personnalisées, Architecture microservices.
